\section{Discussion, Limitations and Conclusion}
\label{sec:discussion}

\subsection{Implications}

The results show why a single comparison between two published factor
implementations is an inadequate diagnosis. A difference can arise from a
controlled configuration change, from the signal and data environment that
remain outside the experiment, from a mismatch between the engine and the
paper's requested method, or from an estimand that is not comparable in the
first place. The factor-specific results range from configuration-sensitive
magnitude despite stable sign and significance (AssetGrowth), to a
configuration-dependent reproduction (MeanRankRevGrowth), to an estimand
mismatch (ShareVol). A paper coefficient and a portfolio-sort spread should
not enter the same replication identity.

The workflow records a paper interpretation in a source-backed
implementation specification and generates the signal formula only after that
specification has been approved. Human review and the fixed configuration menu
govern the empirical choices. Holding the signal fixed permits the same
interpretation to be evaluated under several configurations without allowing
the model to alter the strategy between experiments. The practical gain is a
consistent evidence standard and a versioned record across papers, rather than
a substitute for the researcher's empirical judgment.

The controlled factorial results provide the paper's strongest evidence.
With a complete factorial design, Shapley values allocate the within-engine
configuration difference exactly. For larger contrasts, one-at-a-time (OAT)
ablations provide sensitivity evidence but cannot identify interactions.
External comparisons and the three-term identity remain observational because
they absorb differences in external code, data vintage, sample window, and
estimator.

Robustness in this study does not mean that a factor is invariant across all
plausible implementations. Rather, the controlled results show different
patterns: AssetGrowth remains positive and significant across the baseline,
C\&Z, and standardized-HXZ configurations, while its return magnitude and
$t$-statistic change materially; MeanRankRevGrowth weakens under the
standardized-HXZ configuration; and ShareVol's portfolio-sort tracks remain
insignificant across the configurations examined. The separate-window results
add a temporal qualification: a successful paper-window replication need not
persist after publication. AssetGrowth's baseline loses significance in its
post-publication window, whereas MeanRankRevGrowth remains significant in its
own post-publication window. These robustness results are conditional on the
chosen configurations and windows. The post-publication comparisons are
descriptive and do not identify a causal effect of publication.

\subsection{Limitations}

\subsubsection{Agent and specification limits}
\begin{itemize}
  \item \textbf{No within-agent dispersion study.} Repeated independent
    extractions of each paper have not been run. Consequently, the analysis
    cannot separate variation induced by model sampling from ambiguity in the
    paper or from a genuine extraction error.
  \item \textbf{No cross-model or prompt-sensitivity evaluation.} The study
    does not compare the selected agent model, system prompt, or tool-use
    policy with alternative models or prompting strategies. Its findings
    therefore describe this controlled workflow, not a model-agnostic ranking
    of LLMs or a claim that the same extraction quality will transfer to every
    agent configuration.
  \item \textbf{No blinded evaluation of the diagnosis layer.} Step~8's
    proposal is structurally constrained by deterministic evidence keys and
    identification labels, but this study does not yet compare its claim
    selection or explanatory usefulness with blinded expert judgments.
  \item \textbf{Finite configuration and estimator menu.} The engine clamps
    off-menu or unspecified choices to documented defaults. This preserves
    reproducibility, but means a paper-faithful baseline can contain a scope
    deviation. It also excludes methods such as Fama--MacBeth from the
    controlled portfolio-sort experiment.
  \item \textbf{No external-signal bridge.} C\&Z's independently coded
    signal has not been executed through this engine. The term labelled
    signal-and-environment therefore combines signal, data-vintage, and
    engine differences rather than identifying a pure signal effect.
\end{itemize}

\subsubsection{Experimental and statistical limits}
\begin{itemize}
  \item \textbf{Factorial budget.} Complete factorial attribution is limited
    to three differing fields. Larger contrasts use OAT and cannot identify
    interactions; ShareVol is the present example.
  \item \textbf{External comparability.} The paper, C\&Z, HXZ, and engine
    endpoints may use different windows, data vintages, standard-error
    conventions, or estimands. Exact arithmetic closure of the three-term
    identity does not remove these differences.
  \item \textbf{$t$-channel domain.} The log-$t$ decomposition is available
    only when both compared mean returns are positive and the required signs
    are consistent. Otherwise the system reports an absolute-$t$ change but
    does not manufacture a channel split.
  \item \textbf{Small, illustrative sample.} Three deliberately selected
    factors cannot estimate how frequently each diagnosis occurs in the
    factor literature.
\end{itemize}

\subsection{Future work}

The immediate next step is repeated extraction per paper, which would turn
the current single-draw audit into a measured within-agent dispersion study.
A bridge that runs an external implementation through the same engine and
data snapshot would further separate signal differences from environmental
ones. Extending the menu to additional data sources, estimator families, and
portfolio constructions would reduce selection and clamping limits. Finally,
scaling the protocol to a pre-specified broader factor sample would test
whether the three observed diagnostic patterns generalize beyond these
illustrative cases.

\subsection{Conclusion}

The central result of this thesis is straightforward: a replication gap is
rarely informative on its own. It becomes informative only after the
underlying implementation choices and the estimand being compared are made
explicit. A common label---successful or failed replication---can otherwise
hide differences in weights, breakpoints, data treatment, or the empirical
object itself.

The cases yield three distinct findings. Asset growth is reproduced on the
paper-aligned raw-return metric, although its monthly premium is sensitive to
the configuration. Sales-growth reversal is reproduced under the C\&Z
configuration, whereas the standardized HXZ reference has the opposite sign.
For share turnover, the original paper estimates a Fama--MacBeth
coefficient and the comparison exercises report portfolio-sort spreads. The
paper-anchored numerical decomposition is therefore unavailable.

The workflow makes implementation choices observable and reviewable before
they enter an empirical conclusion. It preserves a traceable record of the
paper interpretation and the assumptions driving a result, supporting research
review and factor validation before implementation.

The framework does not recover the original authors' code or estimate general
LLM extraction accuracy. It documents the comparisons behind a replication
disagreement and flags cases in which the endpoints cannot be compared. These
records provide a basis for subsequent research or implementation decisions.
